% Options for packages loaded elsewhere
\PassOptionsToPackage{unicode}{hyperref}
\PassOptionsToPackage{hyphens}{url}
\PassOptionsToPackage{space}{xeCJK}
\documentclass[
  chinese,
  11pt,
]{ctexart}
\usepackage{xcolor}
\usepackage[margin=2.2cm]{geometry}
\usepackage{amsmath,amssymb}
\setcounter{secnumdepth}{5}
\usepackage{iftex}
\ifPDFTeX
  \usepackage[T1]{fontenc}
  \usepackage[utf8]{inputenc}
  \usepackage{textcomp} % provide euro and other symbols
\else % if luatex or xetex
  \usepackage{unicode-math} % this also loads fontspec
  \defaultfontfeatures{Scale=MatchLowercase}
  \defaultfontfeatures[\rmfamily]{Ligatures=TeX,Scale=1}
\fi
\usepackage{lmodern}
\ifPDFTeX\else
  % xetex/luatex font selection
  \ifXeTeX
    \usepackage{xeCJK}
    \setCJKmainfont[]{Noto Serif SC}
    \setCJKsansfont[]{Noto Sans SC}
    \setCJKmonofont[]{DejaVu Sans Mono}
  \fi
  \ifLuaTeX
    \usepackage[]{luatexja-fontspec}
    \setmainjfont[]{Noto Serif SC}
  \fi
\fi
% Use upquote if available, for straight quotes in verbatim environments
\IfFileExists{upquote.sty}{\usepackage{upquote}}{}
\IfFileExists{microtype.sty}{% use microtype if available
  \usepackage[]{microtype}
  \UseMicrotypeSet[protrusion]{basicmath} % disable protrusion for tt fonts
}{}
\makeatletter
\@ifundefined{KOMAClassName}{% if non-KOMA class
  \IfFileExists{parskip.sty}{%
    \usepackage{parskip}
  }{% else
    \setlength{\parindent}{0pt}
    \setlength{\parskip}{6pt plus 2pt minus 1pt}}
}{% if KOMA class
  \KOMAoptions{parskip=half}}
\makeatother
\usepackage{longtable,booktabs,array}
\newcounter{none} % for unnumbered tables
\usepackage{calc} % for calculating minipage widths
% Correct order of tables after \paragraph or \subparagraph
\usepackage{etoolbox}
\makeatletter
\patchcmd\longtable{\par}{\if@noskipsec\mbox{}\fi\par}{}{}
\makeatother
% Allow footnotes in longtable head/foot
\IfFileExists{footnotehyper.sty}{\usepackage{footnotehyper}}{\usepackage{footnote}}
\makesavenoteenv{longtable}
\ifLuaTeX
\usepackage[bidi=basic,shorthands=off,provide=*]{babel}
\else
\usepackage[bidi=default,shorthands=off,provide=*]{babel}
\fi
\ifLuaTeX
  \usepackage{selnolig} % disable illegal ligatures
\fi
\setlength{\emergencystretch}{3em} % prevent overfull lines
\providecommand{\tightlist}{%
  \setlength{\itemsep}{0pt}\setlength{\parskip}{0pt}}
\usepackage{bookmark}
\IfFileExists{xurl.sty}{\usepackage{xurl}}{} % add URL line breaks if available
\urlstyle{same}
\hypersetup{
  pdftitle={课题02-Transformer骨架},
  pdflang={zh-CN},
  hidelinks,
  pdfcreator={LaTeX via pandoc}}

\title{课题02-Transformer骨架}
\author{}
\date{}

\begin{document}
\maketitle

{
\setcounter{tocdepth}{2}
\tableofcontents
}
\section{课题02 开题简报 · Transformer 骨架与训练循环}\label{ux8bfeux989802-ux5f00ux9898ux7b80ux62a5-transformer-ux9aa8ux67b6ux4e0eux8badux7ec3ux5faaux73af}

\begin{quote}
模块：\textbf{M0/M1/M3} ｜ 周次：\textbf{W2--W3（09-21→10-04）} ｜ 算力：\textbf{T0 纯 CPU，0 GPU·小时} 唐杰作业对应：\textbf{① 从零写 Transformer} 的后半 ｜ 判分来源：\textbf{CS336 A1 \texttt{tests/test\_model.py}、\texttt{test\_nn\_utils.py}、\texttt{test\_optimizer.py}、\texttt{test\_serialization.py}} 手写：\textbf{R2（Transformer 全部组件）+ R3（训练循环与优化器）+ R4（反向传播验证）}
\end{quote}

\begin{center}\rule{0.5\linewidth}{0.5pt}\end{center}

\subsection{一、研究背景}\label{ux4e00ux7814ux7a76ux80ccux666f}

2017 年《Attention Is All You Need》删掉 RNN，用自注意力 + 前馈层实现 \textbf{O(1) 依赖路径 + 全并行}。 此后所有主流大模型都是它的变体，但\textbf{细节已经换了一轮}：正弦位置编码 → RoPE；LayerNorm → RMSNorm； post-norm → pre-norm；ReLU MLP → SwiGLU；MHA → GQA。\textbf{本课题要求你把现代版写出来，而不是复刻 2017 版。}

\subsection{二、研究问题（一句话）}\label{ux4e8cux7814ux7a76ux95eeux9898ux4e00ux53e5ux8bdd}

\begin{quote}
\textbf{一个''能从零训起来''的现代 Transformer 层，最少需要哪几个组件？每个组件去掉之后，训练会以什么方式失败？}
\end{quote}

\subsection{三、攻关线索}\label{ux4e09ux653bux5173ux7ebfux7d22}

\begin{enumerate}
\def\labelenumi{\arabic{enumi}.}
\tightlist
\item
  \textbf{形状先行}：写下每一步的张量形状（含 batch、序列、头数），\textbf{先能在纸上跑通}，再写代码。
\item
  \textbf{注意力的三个动作}：投影出 Q/K/V → 缩放点积 + 因果掩码 → softmax 加权求和并输出投影。哪一步最容易错？为什么？
\item
  \textbf{位置信息}：RoPE 为什么是''旋转 Q/K''而不是''加到输入上''？（讲义 M3 §三给了恒等式）
\item
  \textbf{归一化与残差}：为什么 pre-norm 比 post-norm 好训？残差提供了什么''通路''？
\item
  \textbf{优化器}：AdamW 里''解耦权重衰减''解耦的是什么？为什么它比 L2 正则更常用？
\item
  \textbf{稳定性}：梯度裁剪在什么情况下救命？fp16 vs bf16 的差别如何影响 loss 曲线？
\end{enumerate}

\subsection{四、里程碑}\label{ux56dbux91ccux7a0bux7891}

\begin{itemize}
\tightlist
\item[$\square$]
  M1 手写全部组件：注意力（含因果 mask 与 RoPE）、RMSNorm、SwiGLU FFN、残差堆叠、初始化
\item[$\square$]
  M2 形状与数值判分通过：\texttt{test\_model.py} / \texttt{test\_nn\_utils.py} 全绿
\item[$\square$]
  M3 手写训练循环 + AdamW + LR 调度 + 梯度裁剪 → \texttt{test\_optimizer.py} / \texttt{test\_serialization.py} 全绿
\item[$\square$]
  M4 端到端：\textbf{CPU 上训一个极小模型（d4--d8），loss 稳定下降}，并保存/加载 checkpoint 续训成功
\item[$\square$]
  M5 \textbf{R4 反向传播验证}：与 PyTorch autograd 的数值梯度对照（相对误差 \textless{} 1e-6）
\item[$\square$]
  M6 \textbf{三组改坏实验}（去 mask / 换位置编码 / 去残差）：每组先写预测，再跑，再记录
\item[$\square$]
  M7 一页报告 + 知识卡 ≥3 张
\end{itemize}

\subsection{五、交付物}\label{ux4e94ux4ea4ux4ed8ux7269}

{\def\LTcaptype{none} % do not increment counter
\begin{longtable}[]{@{}lll@{}}
\toprule\noalign{}
交付物 & 位置 & 验收 \\
\midrule\noalign{}
\endhead
\bottomrule\noalign{}
\endlastfoot
实现 & \texttt{手写/}（你写） & M2/M3 判分全绿 \\
判分输出 & \texttt{判卷记录.md} & \texttt{证据/judge-*.log} 原文 \\
改坏实验记录 & \texttt{证据/改坏实验.md} & 预测→现象→解释 三段齐全 \\
loss 曲线 & \texttt{证据/} & 训练日志 + 图（含坐标轴与步数） \\
报告 & \texttt{报告.md} & 一页四段 \\
\end{longtable}
}

\subsection{六、讲课锚点（讲义 M3 + 对标）}\label{ux516dux8bb2ux8bfeux951aux70b9ux8bb2ux4e49-m3-ux5bf9ux6807}

\begin{itemize}
\tightlist
\item
  \textbf{讲义}：\texttt{讲义/M3-Transformer内部机制.md}（主干）+ \texttt{讲义/M0-数学与机器学习地基.md}（§一 自动微分、§五 排查顺序）
\item
  \textbf{对标}：CS336 lecture 03--04 · rasbt ch3--ch4 · 李沐 d2l ch10--ch11
\item
  \textbf{必备工具链}：\texttt{bash\ 环境/bootstrap\_cpu.sh\ -\/-with-torch}（本课题起需要 torch）
\end{itemize}

\subsection{七、代码级提问示例（讲课中出，先预测后验证）}\label{ux4e03ux4ee3ux7801ux7ea7ux63d0ux95eeux793aux4f8bux8bb2ux8bfeux4e2dux51faux5148ux9884ux6d4bux540eux9a8cux8bc1}

\begin{enumerate}
\def\labelenumi{\arabic{enumi}.}
\tightlist
\item
  把因果 mask 从 softmax \textbf{之前}挪到\textbf{之后}，训练现象会怎样？
\item
  去掉残差后，4 层还能训、12 层完全不动------请用梯度路径解释。
\item
  你把 RoPE 的最大频率底数从 10000 改成 100：短文训练会受影响吗？长上下文外推呢？
\item
  AdamW 的 \texttt{weight\_decay} 若改成''直接加到梯度上''（即 L2），在大学习率下会出现什么差异？
\end{enumerate}

\subsection{八、风险与提醒}\label{ux516bux98ceux9669ux4e0eux63d0ux9192}

\begin{itemize}
\tightlist
\item
  ⚠️ \textbf{CPU 训练极慢}：本课题\textbf{只要能训起来、loss 下降}即可，不追求指标（真训练在课题04）。
\item
  ⚠️ 形状错误是最高频故障：\textbf{先写形状注释再写实现}，可省一半时间。
\item
  ⚠️ 判分器需要 torch；\texttt{bootstrap\ -\/-with-torch} 首次约 555MB（沙箱网络受限时可能失败，届时改用 Kaggle CPU 会话）。
\end{itemize}

\end{document}
